\documentclass[12pt, a4paper]{article}

% ── Codificación y lengua ─────────────────────────────────────────────────────
\usepackage[utf8]{inputenc}
\usepackage[T1]{fontenc}
\usepackage[spanish]{babel}
\setlength{\headheight}{15pt}

% ── Matemáticas ───────────────────────────────────────────────────────────────
\usepackage{amsmath, amssymb, amsthm, mathtools}

% ── Algoritmos ────────────────────────────────────────────────────────────────
\usepackage{algorithm}
\usepackage{algpseudocode}
\algnewcommand\algorithmicforeach{\textbf{para cada}}
\algdef{S}[FOR]{ForEach}[1]{\algorithmicforeach\ #1\ \textbf{hacer}}

% ── Código fuente ─────────────────────────────────────────────────────────────
\usepackage{listings}
\usepackage{xcolor}

\definecolor{codebg}{RGB}{248,249,250}
\definecolor{kwcolor}{RGB}{0,90,158}
\definecolor{cmcolor}{RGB}{80,130,80}
\definecolor{stcolor}{RGB}{160,30,30}
\definecolor{numcolor}{RGB}{130,130,130}
\definecolor{rulecolor}{RGB}{200,200,200}

\lstset{
  language=Python,
  backgroundcolor=\color{codebg},
  basicstyle=\ttfamily\footnotesize,
  keywordstyle=\color{kwcolor}\bfseries,
  commentstyle=\color{cmcolor}\itshape,
  stringstyle=\color{stcolor},
  numbers=left,
  numberstyle=\tiny\color{numcolor},
  frame=single,
  framerule=0.5pt,
  rulecolor=\color{rulecolor},
  breaklines=true,
  showstringspaces=false,
  tabsize=4,
  captionpos=b,
  xleftmargin=1.5em,
  framexleftmargin=1.5em,
  aboveskip=10pt,
  belowskip=6pt
}

% ── Tablas ────────────────────────────────────────────────────────────────────
\usepackage{booktabs}
\usepackage{array}
\usepackage{multirow}

% ── Gráficos y diagramas ──────────────────────────────────────────────────────
\usepackage{tikz}
\usetikzlibrary{arrows.meta, positioning, shapes.geometric, fit, backgrounds, calc}

% ── Diseño de página ──────────────────────────────────────────────────────────
\usepackage{geometry}
\geometry{a4paper, top=2.5cm, bottom=2.5cm, left=3cm, right=2.5cm}
\usepackage{setspace}
\onehalfspacing

% ── Encabezados y pie de página ───────────────────────────────────────────────
\usepackage{fancyhdr}
\pagestyle{fancy}
\fancyhf{}
\fancyhead[L]{\small\textit{Ejection Chains en Grafos Bipartitos — GRASP}}
\fancyhead[R]{\small\thepage}
\renewcommand{\headrulewidth}{0.4pt}

% ── Hipervínculos ─────────────────────────────────────────────────────────────
\usepackage{hyperref}
\hypersetup{
  colorlinks=true,
  linkcolor=kwcolor,
  urlcolor=kwcolor,
  citecolor=kwcolor,
  pdftitle={Ejection Chains en Grafos Bipartitos construidos mediante GRASP},
  pdfauthor={},
}

% ── Entornos matemáticos ──────────────────────────────────────────────────────
\theoremstyle{definition}
\newtheorem{definition}{Definición}[section]
\newtheorem{example}{Ejemplo}[section]

\theoremstyle{plain}
\newtheorem{theorem}{Teorema}[section]
\newtheorem{proposition}{Proposición}[section]
\newtheorem{lemma}{Lema}[section]

\theoremstyle{remark}
\newtheorem{remark}{Observación}[section]
\newtheorem{corollary}{Corolario}[section]

% ── Macros propias ────────────────────────────────────────────────────────────
\newcommand{\GE}{G_{\!E}}
\newcommand{\dom}[1]{D\!\left(#1\right)}
\newcommand{\domk}[2]{D_{#1}\!\left(#2\right)}
\newcommand{\inv}[1]{f^{-1}\!\left(#1\right)}
\newcommand{\Nset}[1]{N\!\left(#1\right)}

% ─────────────────────────────────────────────────────────────────────────────
\begin{document}

% ── Portada ───────────────────────────────────────────────────────────────────
\begin{titlepage}
  \centering
  \vspace*{2cm}
  {\LARGE\bfseries Ejection Chains en Grafos Bipartitos\\[0.4em]
   construidos mediante GRASP\par}
  \vspace{1.5cm}
  {\large\itshape Formulaciones, Dominios y Métodos de Búsqueda\\
   con Conjunto Eyectado Mínimo\par}
  \vspace{3cm}
  \rule{0.6\linewidth}{0.8pt}\\[0.4em]
  {\large Informe Técnico\par}
  \rule{0.6\linewidth}{0.8pt}
  \vfill
  {\large \today\par}
\end{titlepage}

% ── Índice ────────────────────────────────────────────────────────────────────
\tableofcontents
\newpage

% =============================================================================
\section{Introducción y Motivación}
% =============================================================================

Las \textbf{Ejection Chains} (cadenas de eyección) son una técnica de búsqueda
local avanzada introducida por Glover~(1991) en el contexto de la metaheurística
Tabu Search. Su principio central consiste en realizar una secuencia de
reasignaciones encadenadas, donde cada movimiento ``eyecta'' un elemento de su
posición actual y lo desplaza a otra, de manera que el elemento desplazado
queda libre para iniciar el siguiente eslabón de la cadena.

Cuando la solución inicial proviene de un procedimiento \textbf{GRASP} (Greedy
Randomized Adaptive Search Procedure), el grafo bipartito resultante posee una
estructura bien definida que permite explotar sistemáticamente los dominios de
los nodos para construir cadenas cortas y de alta calidad.

Este informe presenta:
\begin{itemize}
  \item la formulación matemática del problema de dominios en grafos bipartitos,
  \item la condición de subconjunto $\domk{k}{x_k} \subseteq \dom{x_0}$,
  \item el modelado de la búsqueda de cadena mínima como un problema de camino
        más corto, y
  \item los algoritmos BFS y A* con heurística admisible, junto con estrategias
        de poda de dominio.
\end{itemize}

% =============================================================================
\section{Preliminares y Notación}
% =============================================================================

\begin{definition}[Grafo Bipartito con Asignación]
  Sea $G = (X \cup Y,\, E)$ un grafo bipartito donde $X$ es el conjunto de
  \emph{agentes} e $Y$ el conjunto de \emph{recursos}.  Una
  \emph{asignación factible} es una función inyectiva $f: X \to Y$ tal que
  $(x, f(x)) \in E$ para todo $x \in X$.
\end{definition}

\begin{definition}[Vecindad y Dominio]
  Para un nodo $x \in X$, su \emph{vecindad} en $G$ es
  $\Nset{x} = \{y \in Y \mid (x,y) \in E\}$.
  El \emph{dominio de reasignación} de $x$ dado el estado actual $S$ es:
  \[
    \dom{x} \;=\; \Nset{x} \setminus \{f(x)\},
  \]
  es decir, los recursos alcanzables por $x$ distintos de su asignación
  vigente.
\end{definition}

\begin{definition}[Nodo Eyectado]
  Un nodo $x_k \in X$ queda \emph{eyectado} en el paso $k$ de la cadena
  cuando el nodo anterior $x_{k-1}$ se reasigna al recurso $y_{k-1} = f(x_k)$,
  privando a $x_k$ de su asignación.
\end{definition}

% =============================================================================
\section{Formulación del Dominio en la Ejection Chain}
% =============================================================================

\subsection{Dominio del nodo inicial}

Sea $x_0$ el nodo que se desea mover (nodo inicial de la cadena).  Su dominio
en el paso $0$ es:

\begin{equation}
  \domk{0}{x_0} \;=\; \dom{x_0} \;=\; \Nset{x_0} \setminus \{f(x_0)\}.
  \label{eq:dom0}
\end{equation}

\subsection{Dominio del nodo eyectado en el paso \texorpdfstring{$k$}{k}}

Cuando $x_{k-1}$ elige $y_{k-1} \in \domk{k-1}{x_{k-1}}$, el nodo
$x_k = \inv{y_{k-1}}$ es eyectado.  Su dominio candidato en el paso $k$ es:

\begin{equation}
  \domk{k}{x_k} \;=\;
  \bigl(\Nset{x_k} \setminus \{f(x_k)\}\bigr)
  \;\setminus\;
  \mathcal{V}_{k-1},
  \label{eq:domk}
\end{equation}

donde $\mathcal{V}_{k-1} = \{y_0, y_1, \ldots, y_{k-2}\}$ es el conjunto de
recursos ya visitados en la cadena (filtro anti-ciclo).

\subsection{Condición de subconjunto del dominio inicial}

La condición central para que la cadena sea \emph{válida} respecto al dominio
inicial establece:

\begin{equation}
  \boxed{\domk{k}{x_k} \;\subseteq\; \domk{0}{x_0}}
  \qquad \forall\, k \geq 1.
  \label{eq:subset}
\end{equation}

Esta condición puede imponerse mediante tres mecanismos complementarios:

\paragraph{Filtro directo.} En cada paso $k$, antes de extender la cadena:
\[
  \domk{k,\,\text{filt}}{x_k}
  \;=\;
  \domk{k}{x_k} \cap \domk{0}{x_0}.
\]

\paragraph{Etiquetado de Glover.} Asignar a cada recurso $y \in Y$ una etiqueta
de nivel:
\[
  \ell(y) =
  \begin{cases}
    0 & \text{si } y \in \domk{0}{x_0}, \\
    k & \text{si } y \text{ fue liberado en el paso } k > 0.
  \end{cases}
\]
Solo se permiten extensiones a recursos con $\ell(y) = 0$.

\paragraph{Condición estructural.} Si el grafo satisface
$\Nset{x_k} \subseteq \Nset{x_0}$ para todo $x_k$ en la cadena
(vecindades anidadas), la condición~\eqref{eq:subset} se cumple
automáticamente.

% =============================================================================
\section{Modelado como Problema de Camino Más Corto}
% =============================================================================

El objetivo de mover $x_0 \to y^*$ con el \textbf{número mínimo de eyecciones}
se reduce a encontrar el camino más corto en el \emph{grafo de eyección}
$\GE$.

\begin{definition}[Grafo de Eyección \texorpdfstring{$\GE$}{GE}]
  \label{def:ge}
  El grafo de eyección $\GE = (V_E, A_E)$ asociado a la solución $S$ se
  define por:
  \begin{align}
    V_E &= \bigl\{(x,\, y) \;\bigm|\; x \in X,\; y \in \dom{x}\bigr\},
    \label{eq:VE}\\[4pt]
    A_E &= \Bigl\{(x_i, y_j) \to (x_k, y_\ell) \;\Bigm|\;
            x_k = \inv{y_j},\;
            y_\ell \in \dom{x_k}\Bigr\},
    \label{eq:AE}
  \end{align}
  con función de costo $c\bigl((x_i,y_j),(x_k,y_\ell)\bigr) = 1$ para
  minimizar el número de nodos eyectados.
\end{definition}

\begin{proposition}
  Un camino de longitud $k$ en $\GE$ desde $(x_0, y^*)$ hasta cualquier
  estado $(x_k, y_{\text{libre}})$ con $y_{\text{libre}} \notin \operatorname{Im}(f)$
  corresponde a una ejection chain que mueve $x_0 \to y^*$ eyectando
  exactamente $k$ nodos intermedios.
\end{proposition}

\subsection{Estructura del estado y reconstrucción}

Un estado $(x_k, y_k)$ en $\GE$ codifica: ``el nodo $x_k$ desea ser asignado
al recurso $y_k$''.  La cadena completa se representa como:

\[
  x_0 \xrightarrow{y_0} x_1 \xrightarrow{y_1} \cdots
  \xrightarrow{y_{k-1}} x_k \xrightarrow{y^*} \varnothing,
\]

donde $y^* \notin \operatorname{Im}(f)$ es un recurso libre.

La \textbf{reconstrucción} de la solución se realiza en orden inverso:
\begin{enumerate}
  \item $f(x_k) \leftarrow y^*$
  \item $f(x_{k-1}) \leftarrow y_{k-1}$ \quad (ahora libre)
  \item $\vdots$
  \item $f(x_0) \leftarrow y_0$
\end{enumerate}

\begin{figure}[H]
\centering
\begin{tikzpicture}[
    node distance=1.4cm and 2.2cm,
    agent/.style={circle, draw=kwcolor, fill=kwcolor!10,
                  minimum size=9mm, font=\small\bfseries},
    resource/.style={rectangle, draw=cmcolor, fill=cmcolor!10,
                     minimum width=1.1cm, minimum height=8mm,
                     font=\small\bfseries, rounded corners=2pt},
    arrow/.style={-{Stealth[length=6pt]}, thick},
    dasharrow/.style={{Stealth[length=5pt]}-{Stealth[length=5pt]},
                       dashed, gray, thick},
  ]
  % Agentes (X)
  \node[agent] (x0) {$x_0$};
  \node[agent, below=of x0] (x1) {$x_1$};
  \node[agent, below=of x1] (x2) {$x_2$};

  % Recursos (Y)
  \node[resource, right=2.4cm of x0] (y0) {$y_0$};
  \node[resource, right=2.4cm of x1] (y1) {$y_1$};
  \node[resource, right=2.4cm of x2] (yf) {$y^*$};

  % Asignaciones actuales (gris punteado)
  \draw[dasharrow] (x0) -- node[above,font=\tiny]{$f(x_0)$} ($(x0)!0.5!(y0)+(0,0)$);
  \draw[dasharrow] (x1) -- node[above,font=\tiny]{$f(x_1)\!=\!y_0$} (y0);
  \draw[dasharrow] (x2) -- node[above,font=\tiny]{$f(x_2)\!=\!y_1$} (y1);

  % Reasignaciones de la cadena
  \draw[arrow, kwcolor] (x0) to[bend left=15]
    node[above,font=\small]{eyecta} (y0);
  \draw[arrow, kwcolor] (x1) to[bend left=15]
    node[above,font=\small]{eyecta} (y1);
  \draw[arrow, kwcolor!60!black] (x2) to[bend left=15]
    node[above,font=\small]{libre} (yf);

  % Leyenda
  \node[font=\tiny, gray, below=0.3cm of yf]
    {$y^* \notin \operatorname{Im}(f)$};
\end{tikzpicture}
\caption{Estructura de una ejection chain de longitud 2: $x_0$ eyecta a $x_1$,
         que a su vez eyecta a $x_2$, el cual se reasigna a un recurso libre.}
\label{fig:chain}
\end{figure}

% =============================================================================
\section{Estrategias de Restricción del Dominio}
% =============================================================================

Para mantener cadenas cortas se aplican tres niveles de poda sobre $\dom{x_k}$:

\subsection{Nivel 1 — Filtro estructural (condición de subconjunto)}

\begin{equation}
  \domk{k,\,(1)}{x_k}
  \;=\;
  \domk{k}{x_k} \cap \domk{0}{x_0}.
  \label{eq:filt1}
\end{equation}

Elimina recursos que no pertenecen al espacio de soluciones
alcanzable desde $x_0$.

\subsection{Nivel 2 — Filtro de calidad (GRASP-guided)}

\begin{equation}
  \domk{k,\,(2)}{x_k}
  \;=\;
  \Bigl\{\, y \in \domk{k,\,(1)}{x_k}
  \;\Bigm|\;
  \Delta f(x_k \to y) \;\geq\; \theta \,\Bigr\},
  \label{eq:filt2}
\end{equation}

donde $\theta$ es el umbral mínimo de variación en la función objetivo y
$\Delta f(x_k \to y) = f_{\text{obj}}(S') - f_{\text{obj}}(S)$ con $S'$
obtenida al reasignar $x_k \to y$.

\subsection{Nivel 3 — Filtro anti-ciclo}

\begin{equation}
  \domk{k,\,(3)}{x_k}
  \;=\;
  \domk{k,\,(2)}{x_k} \;\setminus\; \mathcal{V}_{k-1},
  \qquad
  \mathcal{V}_{k-1} = \{y_0, \ldots, y_{k-2}\}.
  \label{eq:filt3}
\end{equation}

Garantiza que la cadena no revisita recursos ya procesados.

\begin{remark}
  La composición de los tres filtros da el \emph{dominio efectivo}:
  \[
    \domk{k,\,\text{eff}}{x_k}
    \;=\;
    \bigl(\dom{x_k} \cap \dom{x_0}\bigr)
    \;\setminus\;
    \bigl(\mathcal{V}_{k-1} \cup \{f(x_k)\}\bigr).
  \]
\end{remark}

% =============================================================================
\section{Algoritmos de Búsqueda}
% =============================================================================

\subsection{BFS — Búsqueda en Anchura (cadena de longitud mínima garantizada)}

BFS es el algoritmo de referencia cuando el costo de arco es uniforme
($c = 1$), ya que garantiza encontrar la cadena con el menor número de
eyecciones.

\begin{algorithm}[H]
\caption{BFS para Ejection Chain Mínima}
\label{alg:bfs}
\begin{algorithmic}[1]
\Require $f$, $\domk{k,\text{eff}}{\cdot}$, $x_0$, $y^*$
\Ensure Camino mínimo $\pi$ o \texttt{Ninguno}
\State $\text{inicio} \leftarrow (x_0,\; y^*)$
\State $Q \leftarrow \langle (\text{inicio},\; [\text{inicio}]) \rangle$
\State $\text{visitado} \leftarrow \{\text{inicio}\}$
\While{$Q \neq \emptyset$}
  \State $((x,\, y),\, \pi) \leftarrow Q.\text{pop\_izquierda}()$
  \If{$y \notin \operatorname{Im}(f)$}
    \Return $\pi$ \Comment{Recurso libre: cadena válida}
  \EndIf
  \State $x' \leftarrow \inv{y}$
  \ForEach{$y' \in \domk{k,\text{eff}}{x'}$}
    \State $s' \leftarrow (x',\, y')$
    \If{$s' \notin \text{visitado}$}
      \State $\text{visitado} \leftarrow \text{visitado} \cup \{s'\}$
      \State $Q.\text{encolar}((s',\; \pi + [s']))$
    \EndIf
  \EndFor
\EndWhile
\Return \texttt{Ninguno}
\end{algorithmic}
\end{algorithm}

\begin{lstlisting}[caption={Implementación BFS en Python}]
from collections import deque

def bfs_ejection(assignment, eff_domain, source, target):
    """
    assignment  : dict  x -> y
    eff_domain  : callable (x, visited) -> set de recursos candidatos
    source      : nodo inicial x_0
    target      : recurso destino y*
    """
    inv = {v: k for k, v in assignment.items()}
    free = set(assignment.values())      # Im(f)

    start = (source, target)
    queue = deque([(start, [start], {target})])  # estado, camino, visitados
    seen  = {start}

    while queue:
        (node, slot), path, visited = queue.popleft()

        if slot not in free:             # recurso libre -> exito
            return path

        evicted = inv[slot]
        for nxt in eff_domain(evicted, visited):
            state = (evicted, nxt)
            if state not in seen:
                seen.add(state)
                queue.append((state, path + [state], visited | {nxt}))

    return None
\end{lstlisting}

\subsection{A* — Búsqueda Heurística (cadena óptima con menor exploración)}

A* reduce el espacio explorado mediante una función heurística admisible
$h: V_E \to \mathbb{R}_{\geq 0}$.

\begin{definition}[Heurística Admisible para Ejection Chains]
  \label{def:heuristic}
  Sea $Y_{\text{libre}} = Y \setminus \operatorname{Im}(f)$ el conjunto de
  recursos libres.  Se define:
  \[
    h(x,\, y) \;=\;
    \begin{cases}
      0 & \text{si } y \in Y_{\text{libre}},\\[3pt]
      1 & \text{si } \exists\, y' \in \dom{\inv{y}} \cap Y_{\text{libre}},\\[3pt]
      2 & \text{en otro caso (cota conservadora).}
    \end{cases}
  \]
\end{definition}

\begin{theorem}
  La heurística de la Definición~\ref{def:heuristic} es \emph{admisible}
  (no sobreestima el número real de pasos restantes) y, por lo tanto,
  A* con esta heurística devuelve una ejection chain de longitud mínima.
\end{theorem}

\begin{proof}[Esquema]
  Si $y \in Y_{\text{libre}}$, el estado ya es meta: $h = 0 \leq 0$ pasos
  reales.  Si el ocupante de $y$ tiene un recurso libre en su dominio,
  se requiere exactamente 1 eyección más: $h = 1$.  En caso contrario,
  se necesitan al menos 2 pasos adicionales: $h = 2$.
  En ningún caso $h$ supera el número real de eyecciones restantes.\qed
\end{proof}

\begin{algorithm}[H]
\caption{A* para Ejection Chain Mínima}
\label{alg:astar}
\begin{algorithmic}[1]
\Require $f$, $\domk{k,\text{eff}}{\cdot}$, $h$, $x_0$, $y^*$
\Ensure Camino óptimo $\pi$ o \texttt{Ninguno}
\State $s_0 \leftarrow (x_0,\, y^*)$
\State $\text{montículo} \leftarrow [(h(s_0),\, 0,\, s_0,\, [s_0])]$
\State $\text{mejor} \leftarrow \{\}$
\While{montículo $\neq \emptyset$}
  \State $(f,\, g,\, (x, y),\, \pi) \leftarrow \text{extraer\_mín}(\text{montículo})$
  \If{$y \in Y_{\text{libre}}$}
    \Return $\pi$
  \EndIf
  \If{$\text{mejor}[(x,y)] \leq g$}
    \State \textbf{continuar}
  \EndIf
  \State $\text{mejor}[(x,y)] \leftarrow g$
  \State $x' \leftarrow \inv{y}$
  \ForEach{$y' \in \domk{k,\text{eff}}{x'}$}
    \State $g' \leftarrow g + 1$;
           \quad $s' \leftarrow (x',\, y')$;
           \quad $f' \leftarrow g' + h(s')$
    \State $\text{insertar}(\text{montículo},\; (f',\, g',\, s',\, \pi+[s']))$
  \EndFor
\EndWhile
\Return \texttt{Ninguno}
\end{algorithmic}
\end{algorithm}

\begin{lstlisting}[caption={Implementación A* en Python}]
import heapq

def heuristic(slot, inv, domains, free):
    if slot in free:
        return 0
    occ = inv.get(slot)
    if occ and any(s in free for s in domains[occ]):
        return 1
    return 2

def astar_ejection(assignment, eff_domain, source, target):
    inv  = {v: k for k, v in assignment.items()}
    free = set(y for y in assignment.values()
               if y not in assignment)   # recursos sin dueno
    # recalcular correctamente
    all_resources = set(assignment.values())
    # free slots = los que no estan en Im(f)
    # (en la practica se pasan como parametro)

    s0 = (source, target)
    h0 = heuristic(target, inv, eff_domain, free)
    heap = [(h0, 0, s0, [s0])]
    best = {}

    while heap:
        f, g, (node, slot), path = heapq.heappop(heap)

        if slot not in all_resources:    # recurso libre
            return path

        if best.get((node, slot), float('inf')) <= g:
            continue
        best[(node, slot)] = g

        evicted = inv.get(slot)
        if evicted is None:
            continue

        visited = {st for _, st in path}
        for nxt in eff_domain(evicted, visited):
            gn = g + 1
            hn = heuristic(nxt, inv, eff_domain, free)
            heapq.heappush(heap, (gn + hn, gn, (evicted, nxt),
                                  path + [(evicted, nxt)]))
    return None
\end{lstlisting}

% =============================================================================
\section{Comparación de Métodos}
% =============================================================================

\begin{table}[H]
\centering
\caption{Comparación de estrategias de búsqueda para ejection chains.}
\label{tab:comparison}
\renewcommand{\arraystretch}{1.3}
\begin{tabular}{lcccc}
\toprule
\textbf{Estrategia} &
\textbf{¿Mín. eyectados?} &
\textbf{Calidad solución} &
\textbf{Memoria} &
\textbf{Complejidad} \\
\midrule
BFS               & Garantizado  & Media        & $O(b^d)$    & $O(b^d)$          \\
A* admisible      & Garantizado* & Alta         & $O(b^d)$    & $O(b^d)$ reducido \\
DFS con límite    & No garantiza & Variable     & $O(d)$      & $O(b^d)$          \\
Greedy local      & Subóptimo    & Alta (local) & $O(1)$      & $O(d)$            \\
Beam Search       & Aproximado   & Alta         & $O(w)$      & $O(w \cdot d)$    \\
\bottomrule
\multicolumn{5}{l}{\small *Óptimo garantizado cuando la heurística es
admisible.  $b$: factor de ramificación; $d$: profundidad; $w$: ancho del haz.}
\end{tabular}
\end{table}

\begin{remark}
  Para grafos con dominios grandes ($|\dom{x}| \gg 1$) y soluciones GRASP
  densas, se recomienda \textbf{A*} con la heurística de la
  Definición~\ref{def:heuristic} combinado con los tres niveles de poda del
  dominio efectivo.  BFS es preferible cuando la heurística no aporta
  discriminación suficiente.
\end{remark}

% =============================================================================
\section{Procedimiento Completo}
% =============================================================================

El Algoritmo~\ref{alg:full} integra la construcción GRASP con la fase de
mejora por ejection chains de longitud mínima.

\begin{algorithm}[H]
\caption{GRASP + Ejection Chains con Conjunto Eyectado Mínimo}
\label{alg:full}
\begin{algorithmic}[1]
\Require Grafo bipartito $G$, función objetivo $f_{\text{obj}}$,
         umbral $\theta$, profundidad máxima $k_{\max}$, iteraciones $I_{\max}$
\Ensure Mejor solución $S^*$ encontrada
\State $S^* \leftarrow \texttt{Ninguno}$
\For{$i = 1, \ldots, I_{\max}$}
  \State $S \leftarrow \texttt{GRASP\_Construir}(G,\, \alpha)$
  \Comment{Fase constructiva voraz-aleatoria}
  \State $S \leftarrow \texttt{BúsquedaLocal}(S)$
  \Comment{Fase de búsqueda local estándar}
  \ForEach{par $(x_0, y^*)$ candidato en $S$}
    \State Calcular $\domk{0}{x_0}$ según~\eqref{eq:dom0}
    \State Construir $\GE$ con dominio efectivo~\eqref{eq:filt3},
           limitado a $k_{\max}$
    \State $\pi \leftarrow \texttt{A*}(\GE,\; (x_0, y^*))$
    \If{$\pi \neq \texttt{Ninguno}$}
      \State $S' \leftarrow \texttt{AplicarCadena}(S,\, \pi)$
      \If{$f_{\text{obj}}(S') > f_{\text{obj}}(S)$}
        \State $S \leftarrow S'$
      \EndIf
    \EndIf
  \EndFor
  \If{$f_{\text{obj}}(S) > f_{\text{obj}}(S^*)$}
    \State $S^* \leftarrow S$
  \EndIf
\EndFor
\Return $S^*$
\end{algorithmic}
\end{algorithm}

% =============================================================================
\section{Ejemplo Ilustrativo}
% =============================================================================

\begin{example}
  Sea $X = \{x_0, x_1, x_2\}$, $Y = \{a, b, c, d\}$,
  con asignación $f(x_0) = a$, $f(x_1) = b$, $f(x_2) = c$
  y dominio $\Nset{x_i} = Y$ para todo $i$.

  Se desea mover $x_0 \to b$:

  \medskip
  \begin{center}
  \begin{tabular}{lll}
  \toprule
  Paso & Acción & Estado \\
  \midrule
  0 & $\domk{0}{x_0} = \{b, c, d\}$ & $x_0$ quiere $b$ \\
  1 & $b$ está ocupado por $x_1$; eyectar $x_1$ & cadena: $(x_0, b)$ \\
  2 & $\domk{1}{x_1} \cap \domk{0}{x_0} = \{c, d\}$ & \\
  3 & $c$ está ocupado por $x_2$; elegir $d$ (libre) & \\
    & $x_1 \to d$ (libre $\checkmark$) & longitud $= 1$ \\
  \midrule
  \multicolumn{3}{l}{\textit{Reconstrucción:} $x_1 \to d$,\; $x_0 \to b$.} \\
  \bottomrule
  \end{tabular}
  \end{center}

  La cadena tiene longitud 1 (un solo nodo eyectado), que es el óptimo para
  este caso.
\end{example}

% =============================================================================
\section{Recomendaciones de Implementación}
% =============================================================================

\begin{enumerate}
  \item \textbf{Estructura de datos.}  Mantener $f^{-1}$ como diccionario
        inverso actualizable en $O(1)$ para calcular al eyectado en cada paso.

  \item \textbf{Dominio efectivo como generador.}  Implementar
        $\domk{k,\text{eff}}{x_k}$ como un generador Python (lazy evaluation)
        para evitar materializar conjuntos grandes.

  \item \textbf{Límite de profundidad $k_{\max}$.}  Establecer $k_{\max} \leq
        |X|/4$ en la práctica; valores mayores rara vez aportan mejoras
        adicionales frente al costo computacional.

  \item \textbf{Ordenamiento de la cola.}  En BFS/A*, desempatar estados con
        igual $f = g + h$ priorizando $-\Delta f_{\text{obj}}$, lo que favorece
        cadenas simultáneamente cortas y de alta calidad.

  \item \textbf{Caché de dominios.}  Si los dominios no cambian entre
        iteraciones GRASP, almacenar $\domk{0}{x}$ para cada $x$ y
        reutilizarlos.

  \item \textbf{Paralelismo.}  Las búsquedas A* desde diferentes nodos
        iniciales $(x_0, y^*)$ son independientes y pueden ejecutarse en
        paralelo con \texttt{concurrent.futures}.
\end{enumerate}

% =============================================================================
\section{Conclusiones}
% =============================================================================

Se ha presentado una formulación rigurosa para la determinación de dominios en
ejection chains sobre grafos bipartitos construidos mediante GRASP, y se ha
modelado el problema de la cadena de longitud mínima como un problema de camino
más corto en el grafo de eyección $\GE$.

Los principales aportes del documento son:

\begin{itemize}
  \item La condición $\domk{k}{x_k} \subseteq \domk{0}{x_0}$ como mecanismo
        unificador del control de dominio.
  \item Tres niveles de poda (estructural, de calidad y anti-ciclo) que reducen
        el espacio de búsqueda sin comprometer la optimalidad.
  \item BFS como garantía de cadena mínima con costo $O(b^d)$.
  \item A* con heurística admisible como alternativa óptima con menor
        exploración efectiva.
  \item Un procedimiento completo GRASP + Ejection Chains (Algoritmo~3)
        listo para implementar.
\end{itemize}

% =============================================================================
% Referencias
% =============================================================================
\newpage
\begin{thebibliography}{9}

\bibitem{glover1991}
  F.~Glover,
  ``Ejection Chains, Reference Structures and Alternating Path Methods
    for Traveling Salesman Problems,''
  \textit{Discrete Applied Mathematics}, vol.~65, pp.~223--253, 1991.

\bibitem{feo1995}
  T.~A.~Feo y M.~G.~C.~Resende,
  ``Greedy Randomized Adaptive Search Procedures,''
  \textit{Journal of Global Optimization}, vol.~6, pp.~109--133, 1995.

\bibitem{glover2000}
  F.~Glover y M.~Laguna,
  \textit{Tabu Search}.
  Kluwer Academic Publishers, 2000.

\bibitem{hart1968}
  P.~E.~Hart, N.~J.~Nilsson y B.~Raphael,
  ``A Formal Basis for the Heuristic Determination of Minimum Cost Paths,''
  \textit{IEEE Transactions on Systems Science and Cybernetics},
  vol.~SSC-4, no.~2, pp.~100--107, 1968.

\bibitem{resende2016}
  M.~G.~C.~Resende y C.~C.~Ribeiro,
  \textit{Optimization by GRASP}.
  Springer, 2016.

\end{thebibliography}

\end{document}
